\hypertarget{addressing-at-8946-nm-the-pic-router-against-crossed-aods-and-an-slmdmd-modulator}{%
\chapter{Addressing at 894.6 nm: the PIC router against crossed AODs and
an SLM/DMD
modulator}\label{addressing-addressing-at-8946-nm-the-pic-router-against-crossed-aods-and-an-slmdmd-modulator}}

Module 31. Code: \texttt{scripts/m31\_addressing.py}; results
\texttt{results/m31\_patterns.json}, \texttt{results/m31\_leakage.json},
\texttt{results/m31\_frequency.json},
\texttt{results/m31\_patterns.png}. Device figures for the PIC come from
modules 4 and 5 (\texttt{docs/findings.md}, Sec. 5). Values are SIM,
DERIVED or SOURCE as marked, and should be reviewed before they are
quoted.

\hypertarget{1-result}{%
\section{Result}\label{addressing-1-result}}

Addressing is where the PIC route has a structural lead over crossed
AODs. It is not alone there: an SLM/DMD modulator shares most of the
lead.

(i) \textbf{Arbitrary patterns.} Crossed AODs light products of a row
set and a column set, so an arbitrary on/off pattern must be split into
such rectangles, one round each. For random patterns this means close to
one round per row:

\begin{itemize}
\tightlist
\item
  7.8 to 16 rounds on a 16 x 16 block (densities 0.05 to 0.5);
\item
  31 to 40 on 40 x 40. With 1 us pulses and 4.5 to 7.5 us AOD access
  times this is 84 to 340 us per pattern. A router with one switch per
  site, or an SLM/DMD modulator, needs one round: 1.2 us for the PIC (26
  ns rise plus 0.19 us of shaped settling) and 24 us for a 43 kHz DMD.
\end{itemize}

(ii) \textbf{One optical frequency.}

\begin{itemize}
\tightlist
\item
  \textbf{AOD.} A crossed AOD tags each site with its own RF offset,
  0.13 MHz per site for a 20 MHz, 150-spot deflector. For a light-shift
  (Z) gate this is harmless. For coherent two-photon addressing against
  a global beam the relative phase walks by 0.84 rad per site separation
  per microsecond, so it must be tracked per site; 40 sites apart, 1 us
  costs 34 rad.
\item
  \textbf{PIC router and DMD.} Both keep one frequency for all sites.
\end{itemize}

(iii) \textbf{Idle-site light.}

\begin{itemize}
\tightlist
\item
  \textbf{PIC, measured.} The measured AlN-on-SiN switch reaches 71.4 dB
  mean extinction and -68 dB on-chip crosstalk. After free-space
  delivery the crosstalk is -50.8 dB, so the delivery optics, not the
  switch, set the floor (\texttt{docs/findings.md}, Sec. 2(iv)).
\item
  \textbf{PIC, simulated.} The cascaded switch of module 4 exceeds 88 dB
  (5th percentile), capped at 70 dB by the assumed stray-light floor.
\item
  \textbf{SLM/DMD.} The demonstrated modulator reaches 46 dB contrast
  and -44 dB nearest-neighbour crosstalk.
\item
  \textbf{AOD.} With k tones on an equally spaced grid, every
  third-order product 2 f\_i - f\_j lands on a site of the same line.
  With 8 tones on a 40-site line, 17 idle sites (54\%) receive such
  light at the intermodulation level, for which we found no sourced
  figure {[}VALUE?{]}.
\end{itemize}

Hence the defensible "better than conventional" claim is about
addressing, not transport:

\begin{itemize}
\tightlist
\item
  arbitrary patterns in one round;
\item
  microsecond switching;
\item
  one frequency;
\item
  extinction above 70 dB at the chip.
\end{itemize}

Its competitor is the SLM/DMD modulator rather than the AOD. It matches
the single round and the single frequency at more than 10,000 sites, but
is about 20 times slower (43 kHz). Its contrast is 46 dB against 71 dB
extinction at the PIC chip, but after free-space delivery the
nearest-neighbour crosstalk differs by only 7 dB (-44 against -51 dB).

\hypertarget{2-arbitrary-patterns}{%
\section{Arbitrary patterns}\label{addressing-2-arbitrary-patterns}}

The number of crossed-AOD rounds is the size of an exact partition of
the pattern\textquotesingle s ones into combinatorial rectangles (each
site lit once). We bound it below by the real rank of the 0/1 matrix and
above by the best of three: a greedy partition, the number of distinct
non-empty rows, and the number of distinct non-empty columns. For random
patterns the bounds nearly coincide:

{\small
\begin{longtable}[]{@{}lllll@{}}
\toprule\noalign{}
block & density 0.05 & 0.10 & 0.25 & 0.50 \\
\midrule\noalign{}
\endhead
\bottomrule\noalign{}
\endlastfoot
4 x 4 & 0.8 & 1.1 & 2.4 & 3.2 \\
8 x 8 & 2.5 & 3.9 & 6.3 & 7.5 to 7.8 \\
16 x 16 & 7.8 & 11.4 & 15.3 to 15.6 & 16.0 \\
24 x 24 & 13.4 to 13.5 & 20.0 to 20.3 & 24.0 & 24.0 \\
40 x 40 & 30.9 to 31.3 & 38.7 to 38.8 & 40.0 & 40.0 \\
\end{longtable}
}

Time per arbitrary pattern on a 16 x 16 block at density 0.25:

{\small
\begin{longtable}[]{@{}llll@{}}
\toprule\noalign{}
pulse & crossed AOD & PIC router & SLM/DMD \\
\midrule\noalign{}
\endhead
\bottomrule\noalign{}
\endlastfoot
0.1 us & 70 to 118 us & 0.32 us & 23 us \\
1 us & 84 to 132 us & 1.22 us & 24 us \\
10 us & 222 to 272 us & 10.2 us & 33 us \\
\end{longtable}
}

Structured patterns (whole rows, checkerboards) need far fewer AOD
rounds; these numbers are for independent sites.

\hypertarget{3-leakage-frequency-channels}{%
\section{Leakage, frequency,
channels}\label{addressing-3-leakage-frequency-channels}}

\textbf{Third-order products} on one AOD line (fraction of idle sites
that receive them):

{\small
\begin{longtable}[]{@{}lllll@{}}
\toprule\noalign{}
sites per line & 2 tones & 4 & 8 & 16 \\
\midrule\noalign{}
\endhead
\bottomrule\noalign{}
\endlastfoot
16 & 6\% & 36\% & 93\% & - \\
40 & 2\% & 15\% & 54\% & 98\% \\
\end{longtable}
}

\textbf{Frequency tags} (0.133 MHz per site): the relative phase after a
pulse is 0.08 / 0.84 / 8.4 rad for adjacent sites at 0.1 / 1 / 10 us,
and 3.4 / 34 / 335 rad for sites 40 apart.

\textbf{Channels and power.}

\begin{itemize}
\tightlist
\item
  PIC: one cascaded switch (two modulators) per site plus slow router
  trims; 16 sites at 20 uW need about 1 mW of input (module 4).
\item
  AOD: 2 RF channels per pair.
\item
  DMD: one device with digital control.
\end{itemize}

The PIC needs as many fast channels as sites; at large site counts that
is the cost of this lead.

\hypertarget{4-limits}{%
\section{Limits}\label{addressing-4-limits}}

\begin{itemize}
\tightlist
\item
  \textbf{Patterns.} Random patterns with independent sites are taken as
  the workload.
\item
  \textbf{AOD intermodulation.} The level is not sourced; we report
  where the products land, not how bright they are.
\item
  \textbf{PIC settling.} The settling time (0.19 us) rests on an assumed
  mechanical mode (20 MHz, Q 10) and ZVD shaping (module 5).
\item
  \textbf{SLM/DMD figures.} These come from one demonstration, read in
  abstract only.
\item
  \textbf{Wavelength.} The measured PIC figures are at 795 nm, not 894.6
  nm.
\end{itemize}

\hypertarget{verification-notes}{%
\section{Verification notes}\label{addressing-verification-notes}}

{\footnotesize
\begin{longtable}[]{@{}>{\raggedright\arraybackslash}p{0.341\dimexpr(\linewidth-8\tabcolsep)\relax}>{\raggedright\arraybackslash}p{0.125\dimexpr(\linewidth-8\tabcolsep)\relax}>{\raggedright\arraybackslash}p{0.108\dimexpr(\linewidth-8\tabcolsep)\relax}>{\raggedright\arraybackslash}p{0.426\dimexpr(\linewidth-8\tabcolsep)\relax}@{}}
\toprule\noalign{}
value & location & type & pointer \\
\midrule\noalign{}
\endhead
\bottomrule\noalign{}
\endlastfoot
rounds per pattern, times & Sec. 1, 2 & SIM &
results/m31\_patterns.json \\
third-order product coverage & Sec. 1, 3 & SIM &
results/m31\_leakage.json \\
frequency tags and phases & Sec. 1, 3 & DERIVED &
results/m31\_frequency.json \\
71.4 dB, -68 dB, -50.8 dB, 26 ns, 8 channels (795 nm) & Sec. 1 & SOURCE
& Zhao et al., arXiv:2508.09920, abstract read 2026-11-15 \\
46 dB, -44 dB, 43 kHz, \textgreater{} 10,000 spots & Sec. 1 & SOURCE &
Zhang et al., arXiv:2310.08539 (Optica 2024), abstract read
2026-11-15 \\
AOD access 4.5 / 7.5 us; 20 MHz band, 150 spots & Sec. 1, 3 & SOURCE &
45035-x-6.5DEG-1.06 (sintec.sg/ao/497.html, read 2026-11-14) \\
0.19 us settling (ZVD, 5\% mode error) & Sec. 1 & SIM & m5 (assumed f0
20 MHz, Q 10) \\
\textgreater= 88 dB cascade, 70 dB stray floor & Sec. 1 & SIM / CHOICE &
m4 (docs/findings.md Sec. 5) \\
AODs limited to row or column patterns, intermodulation limits contrast
& Sec. 1 & SEARCH SUMMARY & Radnaev et al., arXiv:2408.08288
(introduction not read) \\
\end{longtable}
}

Open placeholders: {[}VALUE?{]} AOD third-order intermodulation level at
5 to 20 tones; {[}REF?{]} PIC switch performance at 894.6 nm.
